% BasicSynopsis.tex -- mathematical synopsis of the foundational (Basic) layer of lean-DLMF.
% Compile with LuaLaTeX: latexmk -lualatex BasicSynopsis.tex
\documentclass[11pt,a4paper]{article}
\usepackage{leanreport}
\hypersetup{pdftitle={lean-DLMF: mathematical synopsis of the Basic layer, 2026-11-02}}

\newcommand{\dlmf}{\repofile{lean-DLMF}}
\newcommand{\F}{\mathbf{F}}

\begin{document}

\title{The Basic Layer of \texttt{lean-DLMF}:\\a Mathematical Synopsis}
\author{Bastiaan J. Braams}
\date{2026-11-02}
\maketitle

\begin{abstract}
  The foundational layer of \texttt{lean-DLMF} consists of every \file{Basic.lean} and every
  file below a \file{Basic/} directory, with \module{DLMF.Basic} as its entry point. It is
  meant as the starting point for bringing the special functions of the NIST Digital Library
  of Mathematical Functions (DLMF) into TauCeti and later Mathlib. This synopsis describes,
  chapter by chapter, which functions are defined there, in which normalization and on which
  domain, and which basic properties are proved. Everything in the layer is fully proved, and
  its imports are confined to Mathlib, TauCeti, \module{DLMFMathlib} and other foundational
  files.
\end{abstract}

\section{Overview}

On 2026-11-02 the layer has 235 files with 49{,}068 lines, 1{,}559 definitions and
2{,}000 theorems; 1{,}618 of its declarations carry 1{,}769 DLMF tags, covering 1{,}324
distinct DLMF items. No declaration in the layer is admitted. For each central
function or family of a DLMF chapter it contains:
\begin{itemize}
  \item the definition, with a docstring that states the domain, branch and normalization;
  \item unfolding lemmas and values at the origin, at integer parameters and at exceptional
    parameters;
  \item symmetries, parity and behaviour under conjugation;
  \item continuity, differentiability or analyticity on the stated domain;
  \item the defining differential equation, recurrence or functional equation, and the
    Wronskian where DLMF gives one;
  \item the bridge to the corresponding Mathlib or TauCeti function, where one exists.
\end{itemize}
Asymptotics, zeros, integral representations, expansions, inequalities and transformations
are in the theorem layer and are not described here.

\paragraph{Conventions.} Functions are \emph{totalized}, as in Mathlib. A function defined by
a series is the sum of that series where it converges, and is given a junk value elsewhere;
a multivalued function is its principal branch on a cut plane. The theorems state the domain
on which a property holds. Where Lean's totalized statement holds without DLMF's exclusion
of poles or exceptional parameters, the stronger statement is proved. Complex powers are the
principal powers $z^a=\exp(a\log z)$ with $\operatorname{ph}z\in(-\pi,\pi]$, which agrees with
DLMF \S4.2 and with Mathlib's \lean{Complex.cpow}.

\paragraph{Shared framework.} \module{DLMF.Core.Basic} fixes the conventions used across
chapters: improper integrals on half-lines and on the whole line, Cauchy principal values,
the straight-segment integral $\int_a^b f$ and the primitive $\int_0^z f$, the sector filters
$|\operatorname{ph}z|\le\alpha$, $|z|\to\infty$ (also on the logarithmic plane), finite and
infinite continued fractions with their continuants, the spherical Bessel functions, the
$k$th zero of a function in a region ordered by modulus, and successively extracted
asymptotic coefficients.

\section{Methods (DLMF Chapters 1 and 2)}

\paragraph{Algebraic and analytic methods (Ch.~1).} Most of Chapter 1 is upstream; the Basic
files define what Mathlib and TauCeti lack, and record the upstream names for the rest.
\begin{itemize}
  \item Means, tridiagonal matrices and infinite determinants (\S\S1.2--1.3).
  \item Calculus and vector calculus: the few notions of \S\S1.4--1.6 not in Mathlib
    (Mathlib supplies derivatives, $C^n$, convexity, total variation and $L^p$).
  \item Complex variables (\S\S1.9--1.10): the principal phase is \lean{Complex.arg} with
    values in $(-\pi,\pi]$; contours, winding numbers and residues are TauCeti's.
  \item Zeros of polynomials (\S1.11), on Mathlib's \lean{Polynomial}, its discriminant and
    Descartes sign changes.
  \item Continued fractions $b_0+a_1/(b_1+a_2/(b_2+\cdots))$ in DLMF indexing (\S1.12),
    bridged to Mathlib's \lean{GenContFract}.
  \item Linear second-order ODEs (\S1.13): solution predicates, Wronskians, the standard
    forms.
  \item Integral transforms in DLMF normalization (\S1.14): Fourier, cosine, sine, Laplace,
    Mellin, Hilbert and Stieltjes.
  \item Summability (\S1.15): Abel and Ces\`aro means, the Poisson and Fej\'er kernels, and
    the Riemann--Liouville fractional integrals and derivatives.
  \item Distributions (\S\S1.16--1.17) on Mathlib's $\mathcal D'$ and $\mathcal S'$, and delta
    sequences.
  \item Fourier coefficients and Lebesgue constants (\S1.8), and the Sturm--Liouville
    operators, eigenfunctions and spectral notions of \S1.18 (on Mathlib's Hilbert spaces and
    unbounded operators).
\end{itemize}

\paragraph{Asymptotic approximations (Ch.~2).} The order symbols are Mathlib's
\lean{IsBigO}, \lean{IsLittleO} and \lean{IsEquivalent} along a filter. Asymptotic scales and
expansions with respect to a scale are \lean{Asymptotics.IsAsymptoticScale} and
\lean{Asymptotics.HasAsymptoticExpansion} from \module{DLMFMathlib}. The Basic files add the
data used by \S\S2.3--2.11: the integrals and phase functions of the Laplace, saddle-point and
Bleistein methods, the Mellin-transform method (with Mathlib's \lean{mellin} and
\lean{mellinInv}), the Liouville--Green and difference-equation data, the comparison sums of
\S2.10 and the remainder and Stokes functions of \S2.11.

\section{Elementary functions and the Gamma function (DLMF Chapters 4 and 5)}

\paragraph{Elementary functions (Ch.~4).}
\begin{itemize}
  \item Logarithms and powers (\S4.2): the multivalued and the closed principal logarithm
    and power, the path-integral characterization of $\log$, base changes, modulus and phase
    formulas, and the mapping properties of $\exp$.
  \item The cosecant and secant (\S4.14) with zeros, poles, meromorphy, derivatives and the
    Pythagorean identities; the reciprocal hyperbolic functions (\S4.28).
  \item The principal inverse trigonometric and hyperbolic functions on their cut planes
    (\S\S4.23, 4.37): the inverse equations, analyticity on the open cut domains, the branch
    points, agreement with Mathlib's real inverses, oddness, and the symmetry and complement
    relations.
  \item The Gudermannian and its inverse (\S4.23(viii)), defined by integrals, with closed
    forms and the inverse relation.
  \item The Lambert $W$ function (\S4.13). The map $w\mapsto w+\log w$ is shown to be a
    biholomorphism from the slit plane onto the plane with the half-lines
    $\Im\zeta=\pm\pi$, $\Re\zeta\le-1$ removed. Its inverse (Wright's $\omega$ off its cuts)
    gives $W_k(z)=\omega(\log z+2\pi\iu k)$ on the slit plane, and upper-side limits on the
    negative axis; on $[-1/\eu,0)$ the branches $k=0,-1$ are the real inverses of $y\eu^y$.
    The Basic files prove the inverse equations $W_k(z)\eu^{W_k(z)}=z$, special values, the
    real branches, analyticity and derivatives, and define Wright's $\omega$ and the tree
    function $T$. The existence of the normalized branch family is a theorem, proved from
    Mathlib alone.
\end{itemize}

\paragraph{Gamma function (Ch.~5).} $\Gamma$ and $\psi$ are Mathlib's. The Basic files prove
Weierstrass's product for $1/\Gamma$ and define the real digamma function
\lean{DLMF.Gamma.Real.digamma} (with the bridge to Mathlib's \lean{Complex.digamma}, values at the integers, recurrence and reflection), the scaled Gamma
function $\Gamma^*(z)=\Gamma(z)/(\eu^{-z}z^z(2\pi/z)^{1/2})$ (5.11.3), and Barnes' $G$ by its
Weierstrass product (5.17.3), with the functional equation, values at the integers and the
Glaisher--Kinkelin constant.

\section{Functions defined by integrals (DLMF Chapters 6--8)}

\paragraph{Exponential, sine and cosine integrals (Ch.~6).} The complex $E_1$, $\mathrm{Ein}$,
$\mathrm{Si}$, $\mathrm{Ci}$, $\mathrm{Shi}$, $\mathrm{Chi}$, the real $\mathrm{Ei}$ and
logarithmic integral, and the auxiliary functions $f$, $g$, with the generalized exponential
integral $E_p$ defined through the confluent functions. Proved: power series, entireness of
$\mathrm{Ein}$, $\mathrm{Si}$, $\mathrm{Cin}$, symmetries, the segment-integral forms of \S6.2,
derivatives, and the limits at infinity of the real functions.

\paragraph{Error functions (Ch.~7).} $\operatorname{erf}$ and $\operatorname{erfc}$ are defined
through Kummer's function (7.2.1 via 13.6.7); further the Faddeeva function $w$, Dawson's
integral, the complex Fresnel integrals and auxiliary functions, the Goodwin--Staton integral,
the alternative normalizations of \S7.1, the repeated integrals $\mathrm{i}^n\operatorname{erfc}$
(\S7.18), the Voigt functions (\S7.19) and the inverse error functions (7.17.1). Proved:
values, derivatives, entireness, power series and the real restrictions.

\paragraph{Incomplete Gamma and Beta functions (Ch.~8).} $\gamma(a,z)$, $\Gamma(a,z)$ and the
normalized $P$, $Q$, $\gamma^*$ are defined through the confluent functions on the logarithmic
plane, with the real normalized functions in \module{DLMF.Statistical.Basic}. The incomplete
Beta function $B_x(a,b)$ and $I_x(a,b)$ are defined by the integral for complex $a,b$ with
positive real parts, and continued to complex $z$ through the principal-branch regularized
Gauss function (8.17.8); TauCeti's real function is bridged.

\section{Bessel-type functions (DLMF Chapters 9--12)}

The regularized series $\,{}_0\tilde F_1(;c;w)=\sum w^n/(n!\,\Gamma(c+n))$ of Mathlib is entire
in $w$ for every complex $c$, including the nonpositive integers. Its calculus, collected in
\module{DLMF.Bessel.Basic.RegularizedSeries}, underlies Chapters 9 and 10.

\paragraph{Airy functions (Ch.~9).} With $f(z)={}_0\tilde F_1(;2/3;z^3/9)$ and
$g(z)=z\,{}_0\tilde F_1(;4/3;z^3/9)$,
\[
  \operatorname{Ai}(z)=3^{-2/3}f(z)-3^{-4/3}g(z),\qquad
  \operatorname{Bi}(z)=\sqrt3\bigl(3^{-2/3}f(z)+3^{-4/3}g(z)\bigr)
\]
(9.4.1, 9.4.3). Proved: entireness, Airy's equation $w''=zw$, values at the origin, the
Maclaurin series, the rotated solutions $\operatorname{Ai}(z\eu^{\mp2\pi\iu/3})$, the Wronskians
and connection formulas of \S9.2. Also defined: the modulus and phase functions $M,\theta,N,\phi$
of \S9.8, and Scorer's $\operatorname{Hi}$ (by its Maclaurin series, 9.12.17) and
$\operatorname{Gi}=\operatorname{Bi}-\operatorname{Hi}$, which solve $w''=zw\pm1/\pi$.

\paragraph{Bessel functions (Ch.~10).} $J_\nu$ is Mathlib's \lean{Complex.besselJ}. The Basic
files define $Y_\nu$, $H^{(1)}_\nu$, $H^{(2)}_\nu$, $I_\nu$, $K_\nu$, with integer orders as
limits in the order, and prove Bessel's equation and the modified equation at every order, the
derivatives and basic recurrences, continuity in the order and the Wronskians. On the real
line $I_\nu$ is identified with TauCeti's \lean{Real.besselI}
(\lean{DLMF.Bessel.besselI_ofReal}).

\paragraph{Struve and related functions (Ch.~11).} The Struve functions $\mathbf H_\nu$,
$\mathbf L_\nu$ (and $\mathbf K_\nu$, $\mathbf M_\nu$) by their power series, with the
inhomogeneous Bessel equations (\S11.2); the Lommel functions (\S11.9); the Anger and Weber
functions $\mathbf J_\nu(z)=\pi^{-1}\int_0^\pi\cos(\nu\theta-z\sin\theta)\,\dd\theta$,
$\mathbf E_\nu$ (11.10.1--2) with their calculus.

\paragraph{Parabolic cylinder functions (Ch.~12).} Weber's $U(a,z)$ and $V(a,z)$, Whittaker's
$D_\nu(z)$ and the even and odd solutions $u_1$, $u_2$ of $w''=(z^2/4+a)w$; the cylinder
equations of \S\S12.2, 12.15, the second solution $\bar U$ and the modulus and phase on the real
axis (12.2.22--23). Proved: derivatives, Weber's equation, initial values, uniqueness for given
initial data, the recurrences of $U$, the Hermite specializations, the Wronskians of $U$ with
its rotations and the connection formula 12.2.17; $V$ is entire in $a$ and $z$.

\section{Confluent hypergeometric, Legendre and hypergeometric functions (DLMF Chapters 13--16)}

\paragraph{Confluent hypergeometric functions (Ch.~13).} Kummer's $M(a,c,z)={}_1F_1(a;c;z)$ is
a totalized power series (13.2.2), related to Mathlib's regularized series; Tricomi's $U(a,b,z)$
is defined on the principal branch through the connection formula 13.2.42. Proved: infinite
radius of convergence, Kummer's equation, uniqueness of entire solutions, Kummer's
transformation, the derivative, the contiguous relations of \S13.3(i), conjugation, and the
value of $U$ at zero first parameter. Whittaker's $M_{\kappa,\mu}$, $W_{\kappa,\mu}$ are defined
from $M$ and $U$, with Whittaker's equation, and the coefficients of the Whittaker--Buchholz
expansion (\S13.24) by a generating function.

\paragraph{Legendre functions (Ch.~14).} $P_\nu(z)=F(-\nu,\nu+1;1;(1-z)/2)$, $Q_\nu$, the
Ferrers function $\mathsf P^\mu_\nu$, the associated $P^\mu_\nu$, $Q^\mu_\nu$,
$\boldsymbol Q^\mu_\nu$, $\mathsf Q^\mu_\nu$ and the conical functions, all in the
hypergeometric normalization of \S14.3 through the continued Gauss function, so that $P_\nu$ is
analytic on $\CC\setminus(-\infty,-1]$ for every complex $\nu$. Proved: $P_\nu(1)=1$,
derivatives and Legendre's equation for $P_\nu$ and $\mathsf P^\mu_\nu$. The spherical and
surface harmonics of \S14.30 are defined as functions of angles.

\paragraph{Hypergeometric function (Ch.~15).} ${}_2F_1$ is Mathlib's \lean{ordinaryHypergeometric}
on the unit disk, and the regularized $\F(a,b;c;z)$ is Mathlib's
\lean{regularizedGaussHGFun}. The Basic files contain Pochhammer identities and the Gauss
coefficients; the power series on the disk, the terminating case and absolute convergence on the
circle; analyticity of $\F$ in its parameters for $|z|<1$ (used with the identity principle to
remove parameter restrictions); the derivative
$\tfrac{\dd}{\dd z}F=\tfrac{ab}{c}F(a+1,b+1;c+1;z)$, the hypergeometric equation and its standard
local solutions at $0$, $1$, $\infty$ with their Wronskians. The continuation to the cut plane
$\CC\setminus[1,\infty)$ (\lean{regularizedGaussHG}) is built from Euler's integral and contiguous
shifts and proved to be the analytic continuation of Mathlib's function. Appell's $F_4$
coefficients and series are defined here.

\paragraph{Generalized hypergeometric functions (Ch.~16).} Parameter multisets, coefficients,
finite and terminating series ${}_pF_q$, the normalized principal continuation of
${}_{q+1}F_q$ (defined by a choice from its characterizing properties; its existence and
analyticity on the cut plane, for all complex parameters, are proved here by induction on $q$
through Euler's integral 16.5.2, together with the termwise calculus, the differential
equation and the Euler integrals of the series), bilateral series and logarithmic Barnes integrals, the series facts of
\S16.2, and Meijer's $G$ function with its contours (\S16.17). The Appell and Horn double series
are totalized, with their convergence domains and differential operators; their Euler-integral
branches (\S16.15) use principal powers on stated domains.

\section{$q$-Series and orthogonal polynomials (DLMF Chapters 17 and 18)}

\paragraph{$q$-Hypergeometric functions (Ch.~17).} Finite and infinite $q$-Pochhammer symbols,
integer-index shifted factorials, Gaussian $q$-binomial coefficients, $q$-numbers, Jackson's
$q$-derivative (with DLMF's value at the origin) and its iterates, the Jackson integral and
its bilateral form, and the $q$-Gamma function with its functional equation (5.18). The general
unilateral $\,{}_r\phi_s$ and bilateral $\,{}_r\psi_s$ are defined with DLMF's phase factors,
together with the parameter conditions used in the chapter (regularity, termination, balance,
well-poise). The identities of \S17.2 for shifted factorials (negative indices, base inversion,
reversal, index shifts, multiple and root products, parameter limits) are proved, as are the
Euler-type relations and parity of the $q$-exponentials, $q$-sines and $q$-cosines (\S17.3),
and the product rule and $q\to1$ limit of the $q$-derivative. The four $q$-Appell functions are
defined by their double series.

\paragraph{Orthogonal polynomials (Ch.~18).} A sequence of polynomials orthogonal with respect to a
Borel measure on $\RR$ is defined; it forms a basis, every polynomial is square integrable, and
the three-term recurrence and the Christoffel--Darboux formula are proved. Squared norms,
Christoffel--Darboux and Poisson kernels, moments and Hankel determinants are defined with their
basic identities. The classical families (Jacobi, Laguerre, physicists' Hermite, Legendre,
Gegenbauer, Chebyshev $V,W$; Chebyshev $T,U$ are Mathlib's) are defined, with initial values,
explicit sums, endpoint values, parity and the bridges to Mathlib's Hermite, shifted Legendre and
Chebyshev polynomials. Also defined: the Hahn class and Wilson class of the Askey scheme
(\S\S18.19--18.26), their $q$-analogues (\S\S18.27--18.28), and the associated, Bessel and Pollaczek
polynomials (\S\S18.30, 18.34, 18.35).

\section{Elliptic integrals and elliptic functions (DLMF Chapters 19--23)}

\paragraph{Elliptic integrals (Ch.~19).} The real Legendre integrals $F(\phi|m)$, $E(\phi|m)$,
$\Pi$ and their complete forms for real $m$; Legendre's integrals of complex amplitude, modulus
and characteristic as integrals along the line segment from $0$ to $\phi$ with principal square
roots (\S19.2), bridged to the real ones; the complete $K(k)$ and $K'(k)$ for complex modulus;
Bulirsch's integrals. Carlson's symmetric integrals $R_F$, $R_J$, $R_G$, $R_D$, $R_C$ are defined
by their literal principal-power integrals with admissible-domain predicates and Cauchy principal
values (\S\S19.15--19.16); their elementary symmetries and homogeneity are proved from the
integrals. (The comparison with Carlson's $R$-function needs \texttt{lean-codes} and is in the
theorem layer.)

\paragraph{Theta functions (Ch.~20).} $\theta_j(z|\tau)$, $j=1,\dots,4$, in DLMF's convention
(period $\pi$ in $z$, nome $q=\eu^{\iu\pi\tau}$), defined from Mathlib's \lean{jacobiTheta\_2}
and totalized to $0$ for $\Im\tau\le0$. Also the nome forms $\theta_j(z,q)$, theta constants and
lattice data. Proved: the Fourier series 20.2.1--4, quasi-periodicity under $(m+n\tau)\pi$ and the
half-period translations, parity, analyticity in $z$ and $\tau$, the entire quotients
$\theta_1/\sin z$ and $\theta_2/\cos z$, and the radii of convergence of the theta power series
(\S20.6).

\paragraph{Riemann theta functions (Ch.~21).} For a finite index type $\iota$ (genus
$g=|\iota|$), Riemann matrices are complex symmetric matrices with positive definite imaginary
part (the Siegel upper half-space). The theta series $\theta(z|\Omega)$, its scaled form and the
theta functions with real characteristics are lattice sums over $\ZZ^\iota$. Proved: absolute and
locally uniform convergence on $\CC^\iota\times\mathfrak S_\iota$, continuity and holomorphy in each
variable, the relations for characteristics and half-period characteristics, and the reduction
for $g=1$ to Chapter 20.

\paragraph{Jacobian elliptic functions (Ch.~22).} The twelve functions $\mathrm{pq}(z,k)$,
$\mathrm p,\mathrm q\in\{\mathrm s,\mathrm c,\mathrm d,\mathrm n\}$, defined as theta quotients
with nome $q=\exp(-\pi K'/K)$ and argument $\zeta=\pi z/(2K)$, in lattice form (parameter $\tau$)
and in DLMF's modulus form. Proved for all complex $z$, $k$: Glaisher's relations, parity, and the
theta quarter-period shifts that correspond to $K$ and $\iu K'$. The real Jacobi amplitude
$\operatorname{am}(u|m)$ is the inverse of the order isomorphism $\phi\mapsto F(\phi|m)$ for
$m<1$, with $\mathrm{sn}=\sin\circ\operatorname{am}$, $\mathrm{cn}$, $\mathrm{dn}$, their
derivatives, periodicity and special values, and $\operatorname{arcsn}$. The inverse functions of
\S22.15 with principal ranges, and Jacobi's $\mathcal E(x,k)$ and $Z(x|k)$ (\S22.16) are defined.

\paragraph{Weierstrass functions and modular functions (Ch.~23).} DLMF describes a lattice by
half-periods $\omega_1,\omega_3$ with $\Im(\omega_3/\omega_1)>0$; this is bridged to Mathlib's
\lean{PeriodPair} and its $\wp$. Defined: $\zeta$ and $\sigma$ by their lattice series and product
(\S23.2), the lattice roots $e_j$, the constants $\eta_j$, the invariants $g_2,g_3$, and
$\wp,\zeta,\sigma$ of given invariants (\S23.3(i)). Proved: convergence, parity, analyticity, zeros
and poles of $\zeta$ and $\sigma$, and their relations with $\wp$. Also defined: the action of
$\mathrm{SL}(2,\ZZ)$, modular functions, modular forms and cusp forms, the modular function
$\lambda(\tau)$, the Dedekind sum and the eta multiplier; Klein's $J$ is
\lean{ModularForm.kleinJ}.

\section{Number-theoretic functions (DLMF Chapters 24--27)}

\paragraph{Bernoulli and Euler polynomials (Ch.~24).} $B_n$ is Mathlib's \lean{bernoulli} (DLMF's
convention $B_1=-1/2$), $B_n(x)$ is \lean{Polynomial.bernoulli}; the Euler numbers $E_n=2^nE_n(1/2)$
and polynomials are in \module{DLMFMathlib}. The Basic files give complex evaluation, the Appell
property, derivatives, the generating functions, vanishing and signs, explicit sums, integrality of
$E_n$, the periodic functions and the first entries of Tables 24.2.1--2; the higher-order and
N\"orlund polynomials (\S24.16(i)) with their generating function; and the Genocchi and tangent
numbers (\S24.15).

\paragraph{Zeta and related functions (Ch.~25).} $\zeta$ is Mathlib's \lean{riemannZeta}; the
Dirichlet series, residue, Euler product and functional equation are Mathlib theorems. Defined:
the Stieltjes constants, Riemann's $\xi$, the factor $\chi$ of the functional equation, the
Riemann--Siegel $\theta$ and $Z$ (bridged to Mathlib's \lean{hardyZ}). Proved: the Laurent expansion
at $s=1$, the Dirichlet series of the derivatives and the properties of $\xi$. The Hurwitz zeta
function $\zeta(s,z)$ for complex $z$ is the limit of the Euler--Maclaurin sequence, with existence
of the limit, holomorphy in $s$ and the bridge to Mathlib's \lean{hurwitzZeta}. The polylogarithm
$\mathrm{Li}_s(z)$ and Lerch's $\Phi(z,s,a)$ are the analytic continuations of their power series to
$\CC\setminus[1,\infty)$; the Fermi--Dirac and Bose--Einstein integrals are defined. The polylog
series of natural order, the real dilogarithm as an integral (with its derivative and agreement with
$\mathrm{Li}_2$ on $[-1,1]$) and Catalan's constant $G=\beta(2)$ complete the chapter.

\paragraph{Combinatorial analysis (Ch.~26).} Rising and falling factorials, permutations, set and
integer partitions, lattice paths; the Delannoy, Motzkin and Narayana numbers with their
recurrences (Catalan and Schr\"oder numbers are Mathlib's); restricted partition and composition
counts on Mathlib's \lean{Nat.Partition} and \lean{Composition}; plane partitions as finite
down-closed sets, $pp(n)$ and the symmetry classes; permutation statistics, rook and hit numbers
and $q$-multinomials; DLMF's signed Stirling numbers $s(n,k)$ (Mathlib's $S(n,k)$ is used as is).

\paragraph{Functions of number theory (Ch.~27).} Most arithmetic functions are Mathlib's
($\varphi$, $\sigma_k$, $\mu$, $\lambda$, $\Lambda$, $\omega$, $\Omega$, $\pi$); the Basic files
define the others with their formulas, DLMF's notions of multiplicative and completely
multiplicative functions, Dirichlet characters in DLMF's form with Ramanujan sums and Gauss sums
(bridged to \lean{DirichletCharacter} and \lean{gaussSum}), Lambert series, the representation
numbers $r_k(n)$ and Jacobi's theta series, Euler's product $\prod(1-x^m)$, the discriminant and
Ramanujan's $\tau$.

\section{Functions defined by differential equations (DLMF Chapters 28--32)}

\paragraph{Mathieu functions (Ch.~28).} Mathieu's equation $w''+(a-2q\cos2z)w=0$ with complex
$z,a,q$; its solutions form a space of entire functions determined by initial data at one point,
with existence for arbitrary data. The basic solutions $w_{\mathrm I}$, $w_{\mathrm{II}}$ (28.2.5)
are characterized, with their symmetries 28.2.4, parity, Wronskian, half-period relations
28.2.7--13, separate analyticity in $a$ and $q$, and behaviour under $q\mapsto-q$. Floquet's theorem,
the characteristic equation, the connection formula and the Floquet dichotomy are proved. The
characteristic values $a_n(q)$, $b_n(q)$ for real $q$ are ordered roots of the boundary conditions
of Table 28.2.1; $\mathrm{ce}_n$, $\mathrm{se}_n$ are normalized with DLMF's signs and parity
(Table 28.2.2), with second solutions $\mathrm{fe}_n$, $\mathrm{ge}_n$, Fourier coefficients and
truncated continued fractions. Also defined: the noninteger-order eigenvalues $\lambda_\nu(q)$ and
Floquet solutions $\mathrm{me}_\nu$; the modified Mathieu functions, radial functions and joining
factors (\S\S28.20--28.24); the data of the large-argument and Goldstein expansions and of the
integral relations (\S\S28.25--28.28); Hill's equation and the Whittaker--Hill and Ince equations.

\paragraph{Lam\'e functions (Ch.~29).} Lam\'e's equation
$w''+(h-\nu(\nu+1)k^2\mathrm{sn}^2(z,k))w=0$ and its other forms, the Lam\'e wave equation, and
the Heun form; the four Sturm--Liouville problems of Table 29.3.2 on $[0,K]$, the eigenvalues
$a^m_\nu$, $b^m_\nu$ and normalized eigenfunctions $\mathrm{Ec}^m_\nu$, $\mathrm{Es}^m_\nu$; the
recurrence coefficients, continued fractions, Fourier coefficients, stability notions and
asymptotic coefficients. For integer order the Lam\'e polynomials of the eight types of Table 29.12.1
are constructed explicitly (\S29.15(i)), with the inner product of \S29.14.

\paragraph{Spheroidal wave functions (Ch.~30).} The spheroidal equation and its variants, its
transformed forms, the indicial exponents at $z=1$, the Mathieu special case and the factorization
$(1-x^2)^{m/2}g(x)$; the spectrum and eigenvalues, the normalized angular functions of the first
kind, and the coefficients of continued fractions, recurrences and expansions; the Ferrers-based
objects: functions of the second kind, radial functions and Flammer's normalization.

\paragraph{Heun functions (Ch.~31).} Heun's equation
$w''+(\gamma/z+\delta/(z-1)+\epsilon/(z-a))w'+(\alpha\beta z-q)w/(z(z-1)(z-a))=0$ with
$\epsilon=\alpha+\beta+1-\gamma-\delta$, its normal form and its automorphisms (\S31.2).

\paragraph{Painlev\'e transcendents (Ch.~32).} The six equations PI--PVI and $\mathrm P'_{\mathrm{III}}$
in DLMF's parameter normalization; the B\"acklund maps $T^\pm$ of PII; the Hamiltonian systems
(\S32.6) and Lax pairs (\S32.4); the renormalizations, alternative and symmetric forms
(\S32.2(ii)--(v)); the coalescence cascade PII$\to$PI, \dots, PVI$\to$PV (32.2.24--37), proved as
limits; Painlev\'e II in the Beals--Wong normalization $u''=2u^3+xu-c$ (parameter $c=-\alpha$),
bridged to DLMF's; definitions for special solutions, reductions and applications (\S\S32.8--32.15).

\section{Physical and multivariable functions (DLMF Chapters 33--36)}

\paragraph{Coulomb functions (Ch.~33).} $F_\ell(\eta,\rho)$, $G_\ell$, $H^\pm_\ell$, the
normalizing constant $C_\ell(\eta)$, the Coulomb phase shift $\sigma_\ell(\eta)$ and $\theta_\ell$,
defined through Kummer's and Tricomi's functions (\S33.2), with the elementary properties of the
turning point and $C_\ell$; the auxiliary data of \S\S33.4--33.13; and the functions
$f(\epsilon,\ell;r)$, $h$, $s$, $c$ in the variables $r=-\eta\rho$, $\epsilon=1/\eta^2$ (\S33.14).

\paragraph{$3j$, $6j$, $9j$ symbols (Ch.~34).} On rational angular momenta: triads, admissibility,
the $3j$ symbol by Racah's sum (with the Clebsch--Gordan coefficient), the $6j$ symbol by its Racah
sum, the $9j$ symbol as a sum of products of three $6j$ symbols; the triangle conditions and vanishing
outside the admissible range are proved.

\paragraph{Functions of matrix argument (Ch.~35).} The space of real symmetric $m\times m$ matrices,
the positive definite cone, matrix intervals and integrals over them, $\operatorname{etr}$, the Haar
measure on $\mathrm O(m)$, partitions and partitional shifted factorials, zonal polynomials, the
multivariate Gamma and Beta functions, the Laplace transform and convolution, and the Bessel
functions and ${}_pF_q$ of matrix argument.

\paragraph{Integrals with coalescing saddles (Ch.~36).} The canonical cuspoid and umbilic integrals
with their phase functions and rotated contours, catastrophe sets and scaling data; proved:
integrability of the contour integrands, the elementary special values 36.2.12, 36.2.15, 36.2.18 and
the symmetries 36.2.22--27. The auxiliary moments, series coefficients, approximate zero loci and the
stationary-phase sum used in \S\S36.7--36.11 are defined.

\section{Material from other sources}

Three areas come from Andrews--Askey--Roy and have no DLMF counterpart: Selberg's integral (the
integrands, the unit cube, Aomoto's integrals, and the right-hand side of Selberg's formula with its
convergence domain), harmonic polynomials on $\RR^n$ (their surface inner product, dimension and
orthogonality between degrees), and the real normalized incomplete Gamma and Beta functions of the
statistical area.

\end{document}
